\documentclass[a5paper,10pt]{article}

% https://github.com/InDevRus/filippov-solutions

% Нумерация страниц
\pagenumbering{gobble}

% Настройка полей
\usepackage{geometry}
\geometry{left=1cm}
\geometry{right=1cm}
\geometry{top=1cm}
\geometry{bottom=1.5cm}

% Вставка таблицы
\usepackage{array}
\usepackage{tabularx}

% Инструменты выравнивания формул
\usepackage{amsmath}
\usepackage{mathtools}

% Диакритика в математике
\usepackage{amsfonts}

% Вычисления размеров на ходу
\usepackage{calc}

% Удобные записи производных
\usepackage[italic=true]{derivative}

% Настройки сносок
\usepackage{enotez}

% Длинные знаки векторов
\usepackage{anyfontsize}
\usepackage[b]{esvect}

% Вставка картинок
\usepackage{graphicx}

% Вставка красной строки
\usepackage{indentfirst}
\setlength{\parindent}{1em}

% Условия в командах
\usepackage{xifthen}

% Луа-скрипты
\usepackage{luacode}

% Настройка команд отдельно в математическом и текстовом режимах
\usepackage{mathcommand}

% Для повторения LaTeX кода
\usepackage{multido}

% Конфигурация отступов формул
\delimitershortfall-1sp
\usepackage{mleftright}
\mleftright

% Растягивание объектов
\usepackage{relsize}

% Обтекание текстом
\usepackage{wrapfig}
\usepackage{subcaption}
\usepackage{floatflt}

% Поддержка ввода кириллицы
\usepackage[english, russian]{babel}

% Настройка корневой позиции для компилятора
\usepackage{import}
\providecommand*{\ProjectRootPath}{..}

\import{\ProjectRootPath/preambles/macros/}{theorems}

\import{\ProjectRootPath/preambles/fonts}{font_setup}

% Знак градуса
\usepackage{siunitx}

\import{\ProjectRootPath/preambles/macros/}{operators}
\import{\ProjectRootPath/preambles/macros/}{redefinitions}

\import{\ProjectRootPath/preambles/fonts}{letters_setup}
\import{\ProjectRootPath/preambles/fonts/adjustments}{custom_superscripts}

\import{\ProjectRootPath/preambles/custom}{formatting}
\import{\ProjectRootPath/preambles/custom}{frames}
\import{\ProjectRootPath/preambles/custom}{list_setup}

% TODO: перевести команды на современный стиль объявления
\input{../preambles/plotting}

\begin{document}

Имеем уравнение $y\` = \dfrac {4x} {\pow{x}{2} + \pow{y}{2}}$. 
Для удобства дальнейших построений рассмотрим уравнение над обратной функцией 
$x\br{\y}$, 
используя правило производной обратной функции: 
$x\`\br{\y} = \dfrac {\pow{x}{2} + \pow{y}{2}} {4x}$.

При такой постановке задачи имеем уравнение изоклины: 
$\dfrac {\pow{x}{2} + \pow{y}{2}} {4x} = m$, 
где $m = \dvfrac{x}{y} \linebreak = \dfrac {1} {k} \ne 0$.
$\pow{x}{2} - 4mx + \pow{y}{2} = 0$, $\br{\pow{x}{2} - 4mx + 4\pow{m}{2}} + \pow{y}{2} = 4\pow{m}{2}$.
$\br{x - 2m}^2 + \pow{y}{2} = \br{2m}^2$,
так что изоклинами будут окружности с центром $\br{2m; 0}$
и радиусом $\vbr{2m} > 0$.
Формально говоря, на прямой $x = 0$, $y \ne 0$ касательная 
будет горизонтальной (для исходной задачи), а в точке 
$\br{0; 0}$ поведение производной неясно.
На левом рисунке изображены изоклины, на правом -- решения уравнений.

\begin{figure}[!h]
\begin{minipage}{.5\textwidth}
    \centering
    \begin{tikzpicture}
        \pgfplotsset{
            Graph/.style = {
                samples = 300,
                restrict y to domain = -5 : 5,
            },
            DirectionGraph/.style = {
                samples = 20,
                quiver = {scale arrows = 0.075},
                restrict y to domain = -5 : 5,
            },
            DirectionGraph1/.style = {
                DirectionGraph,
                quiver = {
                    u = {dot_x(x, y) / sqrt(dot_x(x, y) ^ 2 + dot_y(x, y) ^ 2)},
                    v = {dot_y(x, y) / sqrt(dot_x(x, y) ^ 2 + dot_y(x, y) ^ 2)}
                }
            },
            DirectionGraph2/.style = {
                DirectionGraph,
                quiver = {
                    u = {-dot_x(x, y) / sqrt(dot_x(x, y) ^ 2 + dot_y(x, y) ^ 2)},
                    v = {-dot_y(x, y) / sqrt(dot_x(x, y) ^ 2 + dot_y(x, y) ^ 2)}
                }
            },
        }
        \begin{axis}[
            trig format plots = rad,
            width = 7.5cm,
            height = 10cm,
            ymin = -4.2,
            ymax = 4.2,
            xmin = -4.2,
            xmax = 4.2, 
            xtick = {-10, ..., 10},
            ytick = {-10, ..., 10},
            minor tick num = 3,
            declare function = {
                f(\m, \x) = 2 * \m + 2 * \m * cos(x);
                g(\m, \x) = 2 * \m * sin(x);
                dot_x(\x, \y) = x * x + y * y;
                dot_y(\x, \y) = 4 * x;
            }]

            \pgfplotsinvokeforeach {-4, -3.875, ..., -0.125, 0.125, 0.25, ..., 4} {
                \draw (2 * #1, 0) [gray, radius = 2 * abs(#1)] circle;
                    
                \addplot[DirectionGraph1, samples = 30 * abs(#1), domain = -pi * (1 - 0.15) : pi * (1 - 0.15)]
                ({f(#1, x)}, {g(#1, x)});
                \addplot[DirectionGraph2, samples = 30 * abs(#1), domain = -pi * (1 - 0.15) : pi * (1 - 0.15)]
                ({f(#1, x)}, {g(#1, x)});
            }
        \end{axis}
    \end{tikzpicture}
\end{minipage}
\begin{minipage}{.5\textwidth}
    \centering
    \begin{tikzpicture}
        \pgfplotsset{
            Graph/.style = {
                opacity = 0.25,
                samples = 500,
                restrict y to domain = -5 : 5,
            },
            PreciseGraph/.style = {
                Graph,
                thick,
                samples = 1500,
                opacity = 1,
            },
            DirectionGraph/.style = {
                opacity = 0.25,
                samples = 20,
                quiver = {scale arrows = 0.075},
                restrict y to domain = -5 : 5,
            },
            DirectionGraph1/.style = {
                DirectionGraph,
                quiver = {
                    u = {dot_x(x, y) / sqrt(dot_x(x, y) ^ 2 + dot_y(x, y) ^ 2)},
                    v = {dot_y(x, y) / sqrt(dot_x(x, y) ^ 2 + dot_y(x, y) ^ 2)}
                }
            },
            DirectionGraph2/.style = {
                DirectionGraph,
                quiver = {
                    u = {-dot_x(x, y) / sqrt(dot_x(x, y) ^ 2 + dot_y(x, y) ^ 2)},
                    v = {-dot_y(x, y) / sqrt(dot_x(x, y) ^ 2 + dot_y(x, y) ^ 2)}
                }
            },
        }
        \begin{axis}[
            trig format plots = rad,
            width = 7.5cm,
            height = 10cm,
            ymin = -4.2,
            ymax = 4.2,
            xmin = -4.2,
            xmax = 4.2, 
            xtick = {-10, ..., 10},
            ytick = {-10, ..., 10},
            minor tick num = 3,
            declare function = {
                f(\m, \x) = 2 * \m + 2 * \m * cos(x);
                g(\m, \x) = 2 * \m * sin(x);
                dot_x(\x, \y) = x * x + y * y;
                dot_y(\x, \y) = 4 * x;
                p(\C, \x) = sqrt(x^2 - (\C + 2 * ln(x^2 + 4))^2);
                q(\C, \x) = \C - 2 + 2 * ln(x^2 + 4);
            }]

            \pgfplotsinvokeforeach {-4, -3.875, ..., -0.125, 0.125, 0.25, ..., 4} {
                \draw (2 * #1, 0) [gray, radius = 2 * abs(#1)] circle;
                    
                \addplot[DirectionGraph1, samples = 30 * abs(#1), domain = -pi * (1 - 0.15) : pi * (1 - 0.15)]
                ({f(#1, x)}, {g(#1, x)});
                \addplot[DirectionGraph2, samples = 30 * abs(#1), domain = -pi * (1 - 0.15) : pi * (1 - 0.15)]
                ({f(#1, x)}, {g(#1, x)});
            }
            
            \foreach \s in {-1, 1}{
                \addplot[PreciseGraph, domain = 2 : 10] ({\s * p(2 + 6 * ln(0.5), x)}, {q(2 + 6 * ln(0.5), x)});
                
                \addplot[PreciseGraph, domain = 0 : 10, olive] ({\s * p(-6.159, x)}, {q(-6.159, x)});
                \addplot[PreciseGraph, olive] coordinates {(-0.15, -4) (0, -4) (0.15, -4)};
                
                \addplot[PreciseGraph, domain = 0 : 10, blue] ({\s * p(-4.2185, x)}, {q(-4.2185, x)});
                \addplot[PreciseGraph, blue] coordinates {(-0.1, -3) (0, -3) (0.1, -3)};
                
                \addplot[PreciseGraph, domain = 0 : 10, orange] ({\s * p(-2.7726, x)}, {q(-2.7726, x)});
                \addplot[PreciseGraph, orange] coordinates {(-0.01, -2) (0, -2) (0.01, -2)};
                
                \addplot[PreciseGraph, domain = 0 : 10, red] ({\s * p(-2.2189, x)}, {q(-2.2189, x)});
                \addplot[PreciseGraph, red] coordinates {(-0.05, -1) (0, -1) (0.05, -1)};
                
                \addplot[PreciseGraph, domain = 0 : 10, violet] ({\s * p(-2.1299, x)}, {q(-2.1299, x)});
                \addplot[PreciseGraph, violet] coordinates {(-0.05, 1) (0, 1) (0.05, 1)};
                
                \addplot[PreciseGraph, domain = 0 : 10, brown] ({\s * p(-1.99147, x)}, {q(-1.99147, x)});
                \addplot[PreciseGraph, brown] coordinates {(-0.075, 2) (0, 2) (0.075, 2)};
                
                \addplot[PreciseGraph, domain = 0 : 10, magenta] ({\s * p(-1.7346, x)}, {q(-1.7346, x)});
                \addplot[PreciseGraph, magenta] coordinates {(-0.125, 3) (0, 3) (0.125, 3)};
            }
        \end{axis}
    \end{tikzpicture}
\end{minipage}
\end{figure}

\begin{remark} 
    Кривые на рисунке справа построены по общему решению 
    исходного уравнения, которое имеет неявную форму 
    $\pow{x}{2} + \br{\y + 2}^2 = C\pow{e}{\elevate[0.3em]{\frac {y}{2}}} - 4$.
\end{remark} 

\end{document}
